\documentclass[10pt,a4paper]{amsart}
\usepackage[margin=19mm]{geometry}
\usepackage{amsmath,amssymb,mathtools}
\usepackage[scr=rsfs,cal=euler]{mathalfa}
\usepackage{xfrac}
\usepackage[unicode,hidelinks,bookmarks=false]{hyperref}
\newcommand{\ie}{i.e.\ }
\newcommand{\cf}{cf.\ }
\newcommand{\resp}{respectively\ }
\newcommand{\Subsection}{\subsection}
\providecommand{\nobreakdash}{\nobreak\hbox{-}}
\setlength{\emergencystretch}{2em}
\multlinegap0pt
\usepackage{bm}
\usepackage{bbm}
\usepackage{dsfont}
\usepackage{xspace}
\usepackage{cancel}
\usepackage{mathtools}
\usepackage{amsthm}
\makeatletter
\@ifundefined{lemma}{\newtheorem{lemma}{Lemma}}{}
\makeatother
\title[On Cloitre's hiccup sequences]{On Cloitre's hiccup sequences}
\author{Robbert Fokkink, Gandhar Joshi}
\date{Source: arXiv:2507.16956v4 (CC BY 4.0)}
\begin{document}
\maketitle

\noindent\emph{Source and licence.} On Cloitre's hiccup sequences. Version arXiv:2507.16956v4 (CC BY 4.0), distributed under the Creative Commons Attribution 4.0 International licence, as recorded in the arXiv licence metadata for this exact version and shown on the abstract page https://arxiv.org/abs/2507.16956v4. Full rights notice: licenses/arxiv-2507.16956-CC-BY-4.0.txt.\par\smallskip

\noindent\textit{Editorial scope.} Counted unit: the Lemma whose source label is \texttt{lem:estimate} in the article (source0.tex lines 1281--1284, source label \texttt{lem:estimate}), delivered together with the article's own complete original proof of it (source0.tex lines 1285--1365, one \texttt{proof} environment of 81 lines). No mathematics is written, repaired or re-derived: statement and proof are the article's own text line for line. Required context retained uncounted: none. Every definition, display and label of those lines is retained unchanged. Omitted from this package and not redistributed: 1--1280, 1366--1485. Nothing inside a retained excerpt is omitted or altered: every retained body is delivered exactly as the article's lines are written, so no omission receipt of any kind accompanies this package. Citations: the retained excerpts carry no citation command at all, so no bibliography entry of the article is transcribed and no reference is redistributed. Reference adaptation: none was needed and none was made; every reference inside the delivered unit resolves inside it, so no unresolved reference or citation is produced. Printed slips kept literal: no printed word, symbol, hypothesis or number of the counted statement or of its proof was altered. Layout and class: the article's own class is \texttt{article} with packages including color, subcaption, amssymb, amsmath, amsthm, amsfonts, amscd, graphicx, tikz, enumitem, mathtools, geometry, bm, hyperref; the wrapper is the lane's pinned \texttt{amsart} editorial wrapper at 10pt on A4, which restates the theorem-like environments the retained text needs and adds the article's own macro definitions that the retained text needs (none), with extra wrapper packages (bm, bbm, dsfont, xspace, cancel, mathtools). The delivered numbering is the unit's own. Assets: no retained span contains a figure environment, a graphics-inclusion command, a PDF-page inclusion, a table or a TikZ picture; no image file is copied into this package, the package declares no asset dependency and no image file is redistributed with it. Licence: version arXiv:2507.16956v4 of this article is distributed under the Creative Commons Attribution 4.0 International licence, as recorded in the arXiv licence metadata for this exact version and shown on the abstract page https://arxiv.org/abs/2507.16956v4. A full rights notice is delivered at licenses/arxiv-2507.16956-CC-BY-4.0.txt. Characters: every retained excerpt is pure ASCII, so the delivered engine, XeLaTeX with its default Latin Modern text font, reports no missing character.
\begin{lemma}\label{lem:estimate}
Let $\psi=-\frac {1}{\phi}$.
For each $n\geq 1$ we have $r_n\in \left(\phi+\frac{1}{n(n+1/2)\sqrt 5},\phi+\frac{1-1/{(5n^2)}}{n^2\sqrt 5}\right)$.
\end{lemma}

\begin{proof}
    Let $J_n=\left(\phi+\frac{1}{n(n+1/2)\sqrt 5},\phi+\frac{1-\psi^{4n}}{n^2\sqrt 5}\right)$. We prove that \[f_n(J_n)\supset J_{n+1}.\] Since $f_n$ is decreasing, we need to verify the two inequalities
    \begin{equation}\label{left}
    f_n\left(\phi+\frac{1}{n(n+1/2)\sqrt 5}\right)>\phi+\frac{1-\psi^{4(n+1)}}{(n+1)^2\sqrt 5},        
    \end{equation}
    and
    \begin{equation}\label{right}
    f_n\left(\phi+\frac{1-\psi^{4n}}{n^2\sqrt 5}\right)<\phi+\frac{1}{(n+1)(n+3/2)\sqrt 5}.
    \end{equation}
    We have that 
\begin{align*}
f_n(\phi + x) &= \frac{n}{n+1} \cdot \frac{1}{\phi + x - 1 - \frac{1}{n\phi}} 
= \frac{n}{n+1} \cdot \frac{1}{x + \frac{n-1}{n\phi}} \\
&= \frac{n^2\phi}{(n+1) n \phi x + n^2 - 1} 
= \frac{\phi}{1 + \frac{\phi x (n+1)}{n} - \frac{1}{n^2}}\\
&=\frac{\phi}{1 - \frac{1-\phi x n(n+1)}{n^2}}.
\end{align*}
    For the first inequality~\ref{left} we have $x=\frac{1}{n(n+1/2)\sqrt 5}$ and therefore we need to verify that
    \[
    \frac{\phi}{1-\frac {\sqrt 5(n+1/2) -\phi(n+1)}{n^2(n+1/2)\sqrt 5}}>\phi+\frac{1-1/(5(n+1)^2)}{(n+1)^2\sqrt 5}.
    \]
    We will prove the sharper inequality
    \[
    \frac{\phi}{1-\frac {\sqrt 5(n+1/2) -\phi(n+1)}{n^2(n+1/2)\sqrt 5}}>\phi+\frac{1}{(n+1)^2\sqrt 5}.
    \]
    This can be rewritten as
    \[
    \frac {\phi(\sqrt 5(n+1/2) -\phi(n+1))}{n^2(n+1/2)\sqrt 5}>\frac{1}{(n+1)^2\sqrt 5}\left(1-\frac {\sqrt 5(n+1/2) -\phi(n+1)}{n^2(n+1/2)\sqrt 5}\right).
    \] 
    Since $\phi(\sqrt 5 -\phi)=1$ this is
    \[
    \frac{n+1-\phi/2}{n+1/2}>\frac{n^2}{(n+1)^2}\left(1-\frac {\sqrt 5(n+1/2) -\phi(n+1)}{n^2(n+1/2)\sqrt 5}\right).
    \]
    Both factors on the right-hand side are less than $1$ and the inequality follows from
    \[
    \frac{n+1-\phi/2}{n+1/2}>\frac{n^2}{(n+1)^2},
    \]
    which reduces to $(n+1-\phi/2)(n+1)^2>n^2(n+1/2)$. It is straightforward to check that this holds
    by comparing the coefficients of the polynomials.
    
    For the
    second inequality in~\ref{right} we have that $x=\frac{1-1/(5n^2)}{n^2\sqrt 5}$ and we need to verify that 
    \[
    \frac{\phi}{1-\frac {n\sqrt 5 -\phi(n+1)(1-1/(5n^2))}{n^3\sqrt 5}}<\phi+\frac{1}{(n+1)(n+3/2)\sqrt 5}.
    \]
    For $n=1$ the left-hand side is $1.397...$ and the right-hand side is $1.707...$. For $n=2$ 
    the left-hand side is $1.605...$ and the right-hand side is $1.660...$. 
    %For $n=3$ we get $1.628...$ and $1.642...$.
    We may assume that $n>2$.
    
    The inequality can be rewritten as
    \[
    \frac{\phi(n\sqrt 5 -\phi(n+1)(1-1/(5n^2)))}{n^3\sqrt 5}<\frac{1}{(n+1)(n+3/2)\sqrt 5}\left(1-\frac {n\sqrt 5 -\phi(n+1)(1-1/(5n^2))}{n^3\sqrt 5}\right),
    \]
    which simplifies to
    \[
    \frac{
    n-\phi^2+\phi^2/{(5n)}+\phi^2/(5n^2)}n<\frac{n^2}{(n+1)(n+3/2)}\left(1-\frac {n/\phi -\phi +\phi/(5n)+\phi/(5n^2)}{n^3\sqrt 5}\right).
    \]
    The left-hand side is \[1-\phi^2\cdot \frac{1}{n}+O\left(\frac 1{n^2}\right)\] and the right-hand side is \[1-\frac 52 \cdot \frac 1n+O\left(\frac 1 {n^2}\right)\] and therefore the inequality holds for sufficiently large $n$. We need to verify that it
    holds for all $n$. On the right-hand side of the equation, the factor 
    \[
    1-\frac {n/\phi -\phi +\phi/(5n)+\phi/(5n^2)}{n^3\sqrt 5}
    \]
    is less than $1$ if $n>2$. If we ignore that factor, we get
    \[
    \frac{
    n-\phi^2+\phi^2/{(5n)}+\phi^2/(5n^2)}n<\frac{n^2}{(n+1)(n+3/2)}.
    \]
Since $\frac{n^2}{(n+1)(n+3/2)}>1-\frac 52 \cdot\frac 1n$ it suffices to show that
    \[
    \frac{
    n-\phi^2+\phi^2/{(5n)}+\phi^2/(5n^2)}n=1-\phi^2\cdot\frac 1n +\frac{\phi^2}{5n^2+5n^3}<1-\frac 52 \cdot\frac 1n,
    \]
    We arrive at
    \[
    \frac{\phi^2}{5n+5n^2}<\frac 1\phi-\frac 12 \text{ for }n>2.
    \]
    It suffices to check this at $n=3$, when the left-hand side is $0.04...$ and the right-hand side is $0.11..$.
    The inequality holds.
\end{proof}

\end{document}
